\section{長方形による求積}\label{節:導入長方形}
速度のグラフが直線なら,その下の部分は三角形なので面積がすぐ分かった.では,グラフが曲線ならどうだろうか.\,$y=x^2$ のグラフの下で,\,$0\le x\le1$ の範囲にある部分の面積を $S$ と書き,速度から位置を求めたときと同じように長方形で求めてみよう.

横の区間を $n$ 等分すると,1つの小区間の幅は $\frac1n$ である.\,$x^2$ はこの範囲で増加するから,小区間の左端の高さで作った長方形は曲線の下に収まり,右端の高さで作った長方形は曲線の上まで覆う.左端の高さを使った和を $L_n$,右端の高さを使った和を $U_n$ とすると,
\begin{align*}
 L_n&=\frac1n\sum_{k=0}^{n-1}\left(\frac{k}{n}\right)^2,\\
 U_n&=\frac1n\sum_{k=1}^{n}\left(\frac{k}{n}\right)^2
\end{align*}
である.添字が異なるのは,左端では $0$ から $\frac{n-1}{n}$ まで,右端では $\frac1n$ から $1$ までの高さを使うためである.図\ref{図:導入長方形}では $n=4$ としている.面積は両者の間にあるから,
\[
 L_n\le S\le U_n
\]
が成り立つ.

\begin{figure}[htbp]
\centering
\begin{tikzpicture}[x=3.7cm,y=3.1cm,>=Stealth]
 \begin{scope}
  \foreach \k in {0,1,2,3}{
   \pgfmathsetmacro{\xl}{\k/4}
   \pgfmathsetmacro{\xr}{(\k+1)/4}
   \pgfmathsetmacro{\hh}{(\k/4)^2}
   \filldraw[fill=blue!16,draw=blue!60] (\xl,0) rectangle (\xr,\hh);
  }
  \draw[->] (-0.08,0)--(1.13,0) node[right]{$x$};
  \draw[->] (0,-0.05)--(0,1.17) node[above]{$y$};
  \draw[thick,domain=0:1,samples=60] plot (\x,{\x*\x});
  \node[below] at (1,0) {$1$};
  \node at (0.5,-0.23) {左端の高さ：$L_4$};
 \end{scope}
 \begin{scope}[xshift=5.4cm]
  \foreach \k in {0,1,2,3}{
   \pgfmathsetmacro{\xl}{\k/4}
   \pgfmathsetmacro{\xr}{(\k+1)/4}
   \pgfmathsetmacro{\hh}{((\k+1)/4)^2}
   \filldraw[fill=red!13,draw=red!60] (\xl,0) rectangle (\xr,\hh);
  }
  \draw[->] (-0.08,0)--(1.13,0) node[right]{$x$};
  \draw[->] (0,-0.05)--(0,1.17) node[above]{$y$};
  \draw[thick,domain=0:1,samples=60] plot (\x,{\x*\x});
  \node[below] at (1,0) {$1$};
  \node at (0.5,-0.23) {右端の高さ：$U_4$};
 \end{scope}
\end{tikzpicture}
\caption{下からの長方形和と上からの長方形和.分割を細かくすると両者の差が小さくなる.}
\label{図:導入長方形}
\end{figure}

まずは実験をしてみよう.\,$n$ をいくつか選んで $L_n,U_n$ を計算すると,次の表のようになる.
\begin{center}
\begin{tabular}{cccc}
\toprule
$n$ & $L_n$ & $U_n$ & $U_n-L_n$\\
\midrule
$1$ & $0$ & $1$ & $1$\\
$2$ & $0.125$ & $0.625$ & $0.5$\\
$4$ & $0.21875$ & $0.46875$ & $0.25$\\
$8$ & $0.2734375$ & $0.3984375$ & $0.125$\\
$100$ & $0.32835$ & $0.33835$ & $0.01$\\
\bottomrule
\end{tabular}
\end{center}
$n$ を大きくすると $L_n$ と $U_n$ の差は縮まり,どちらも $0.333\ldots$ に近づいていくように見える.\,$S=\frac13$ と予想できるだろう.また,差 $U_n-L_n$ はちょうど $\frac1n$ になっている.

先に差の方を確かめよう.\,$U_n$ と $L_n$ は,高さ $\left(\frac1n\right)^2,\ldots,\left(\frac{n-1}n\right)^2$ の項を共通に含む.引き算するとこれらが打ち消し合い,
\begin{align*}
 U_n-L_n
 &=\frac1n\left\{\left(\frac nn\right)^2-\left(\frac0n\right)^2\right\}\\
 &=\frac1n
\end{align*}
だけが残る.図で言えば,右端の高さの長方形を幅 $\frac1n$ だけ右へずらすと左端の高さの長方形に重なり,最後の高さ $1$ の長方形だけがはみ出すのである.

次に値そのものを求めよう.自然数の平方の和の公式
\[
 \sum_{k=1}^{n}k^2=\frac{n(n+1)(2n+1)}6
\]
を使うと,
\begin{align*}
 U_n&=\frac1{n^3}\sum_{k=1}^{n}k^2\\
 &=\frac{n(n+1)(2n+1)}{6n^3}\\
 &=\frac13+\frac{1}{2n}+\frac{1}{6n^2}
\end{align*}
となる.\,$L_n$ は $U_n$ から $\frac1n$ を引けばよいので,
\[
 L_n=\frac13-\frac{1}{2n}+\frac{1}{6n^2}
\]
である.\,$n$ を大きくすれば,\,$L_n$ も $U_n$ も $\frac13$ に近づく.\,$S$ は $n$ によらない値で,常に $L_n\le S\le U_n$ を満たすのだから,\,$S=\frac13$ と考えられる.予想どおりである.ただし,ここでは曲線に囲まれた部分の面積が定まっていることを前提にした.この議論を,面積の定義まで含めて明確にするのが定積分である.

差 $U_n-L_n=\frac1n$ は,未知の $S$ をどれだけの幅に閉じ込めたかを表している.\,$S$ の値を知らなくても,\,$U_n$ を近似値としたときの誤差は $\frac1n$ 以下だと分かるのだ.

\begin{example}{面積の近似}
上の放物線の面積を $10$ 等分した長方形で挟め.さらに,右端の和 $U_{10}$ を面積の近似値としたとき,どの程度大きく見積もる可能性があるか説明せよ.
\end{example}
\paragraphbf{方針}
$L_n,U_n$ の式に $n=10$ を代入して $S$ を挟む.大きく見積もる誤差は,\,$S$ を使わずに $U_{10}-L_{10}$ で上から抑える.
\paragraphbf{解答}
$0.285\le S\le0.385$.大きく見積もる誤差は $0.1$ 以下である.
\paragraphbf{解法}
$n=10$ を代入すると,
\begin{align*}
 L_{10}&=\frac13-\frac1{20}+\frac1{600}\\
 &=\frac{57}{200}=0.285,\\
 U_{10}&=\frac13+\frac1{20}+\frac1{600}\\
 &=\frac{77}{200}=0.385
\end{align*}
である.したがって $0.285\le S\le0.385$ と分かる.\,$U_{10}$ は常に $S$ 以上であり,
\[
 0\le U_{10}-S\le U_{10}-L_{10}=0.1
\]
だから,\,$U_{10}$ を近似値としたときに大きく見積もる誤差は $0.1$ 以下である.この評価には $S$ の値を使っていない.実際の面積 $\frac13$ を知った後なら誤差は $U_{10}-\frac13=\frac{31}{600}$ と計算できるが,未知の値を求める途中で使えるのは前者の評価である.
